\chapter*{Príloha A: Materiály k testovaniu aplikácie}
\addcontentsline{toc}{chapter}{Príloha A}

Nasledujúce výroky slúžia na vyjadrenie Vašich skúseností s používaním
alternatívnych editorov. Pri každom prosím označte mieru svojho súhlasu na škále od 1 do 5, kde:

\begin{itemize}
    \item 1 = rozhodne nesúhlasím
    \item 5 = rozhodne súhlasím
\end{itemize}

\hrule
\vspace{1em}

\noindent
\begin{tabularx}{\linewidth}{@{}>{\raggedright\arraybackslash}m{0.6cm} >{\raggedright\arraybackslash}X >{\raggedleft\arraybackslash}m{2cm}@{}}

    1.  & Myslím si, že by som alternatívne editory chcel/a používať často.
        & \makebox[1cm][l]{\dotfill}                                                                                            \\[1em]

    2.  & Alternatívne editory sa mi zdali zbytočne zložité.
        & \makebox[1cm][l]{\dotfill}                                                                                            \\[1em]

    3.  & Myslím si, že alternatívne editory boli jednoduché na použitie.
        & \makebox[1cm][l]{\dotfill}                                                                                            \\[1em]

    4.  & Myslím si, že by som potreboval/a pomoc technicky zdatnej osoby, aby som vedel/a tieto alternatívne editory používať.
        & \makebox[1cm][l]{\dotfill}                                                                                            \\[1em]

    5.  & Jednotlivé funkcie alternatívnych editorov boli podľa mňa dobre integrované.
        & \makebox[1cm][l]{\dotfill}                                                                                            \\[1em]

    6.  & Myslím si, že v alternatívnych editoroch bolo príliš veľa nekonzistentností.
        & \makebox[1cm][l]{\dotfill}                                                                                            \\[1em]

    7.  & Predpokladám, že väčšina ľudí by sa naučila používať tieto alternatívne editory veľmi rýchlo.
        & \makebox[1cm][l]{\dotfill}                                                                                            \\[1em]

    8.  & Tieto alternatívne editory sa mi zdali veľmi ťažkopádne na používanie.
        & \makebox[1cm][l]{\dotfill}                                                                                            \\[1em]

    9.  & Pri používaní alternatívnych editorov som sa cítil/a veľmi sebaisto.
        & \makebox[1cm][l]{\dotfill}                                                                                            \\[1em]

    10. & Musel/a som sa naučiť veľa vecí, kým som vedel/a s týmito alternatívnymi editormi začať pracovať.
        & \makebox[1cm][l]{\dotfill}                                                                                            \\[1em]
\end{tabularx}

\vspace{1em}

\hrule

\vspace{1em}

\noindent
Sem môžete uviesť konkrétne pripomienky k alternatívnym editorom a prípadne aj návrhy na ich zlepšenie.

\noindent\dotfill

\begin{figure}
    \centerline{\includegraphics[width=0.98\textwidth]{\imgdir/test-structure-whole.png}}
    %popis obrazku
    \caption[Štruktúra z cvičenia zameraného na testovanie alternatívnych editorov.]{Štruktúra z cvičenia zameraného na testovanie alternatívnych editorov.}
    %id obrazku, pomocou ktoreho sa budeme na obrazok odvolavat
    \label{obr:test-struct}
\end{figure}